\documentclass[11pt]{article}
\usepackage[margin=1in]{geometry}
\usepackage{amsmath,amssymb}
\usepackage{booktabs}
\usepackage{url}
\title{Rewriting the manuscript around one narrative, with the PCBF guarantee
under the estimator's certified uncertainty}
\date{Prepared October 11, 2026}
\begin{document}
\maketitle

\section{Question}
The user asked for \path{paper/manuscript.tex} to be rewritten so that it reads
as natural, human-written prose and follows one narrative: lane keeping and
collision avoidance should be one system; existing designs assume known states
and give no assurance that a safe action still exists a moment later; Sharma
et al.\ (2026) showed that tracking from a maneuvering car needs an
ego-motion-aware relative motion model but assumed constant ego acceleration
and yaw rate and reported no usable error bound; Huang et al.\ (2025) turned a
finite-horizon problem into a barrier function but assumed an exact state and
a time-invariant problem; our estimator and controller are designed for each
other. The user named the remaining gap, that Huang's guarantees are
established only with exact information so the estimator's certificate does
not enter the safety proof, and asked that the mathematics be developed in
the paper for the PCBF guarantee under the estimator's certified uncertainty
and under a forecast that is revised at every sample.

\section{What was read}
The whole manuscript of the previous revision (\texttt{1000635}, synchronized
with \texttt{ea88059}); the two reference papers (Sharma, Alai and Rajamani,
Transp.\ Res.\ C 182, 2026: models (2), (17), Assumption~1, Theorems~1--2;
Huang, Wang, Margellos and Goulart, ECC 2025: problem (8), the shifted
solution (10), inequality (11), Lemmas III.1 and III.5, Theorems~1--4);
\path{controller/OBSERVER_ROBUST_PCBF_THEORY.tex},
\path{controller/TERMINAL_SAFE_SET.md} (Sections 6 and 9),
\path{controller/PCBF_CLF_ARCHITECTURE.md}, \path{controller/NOMINAL_CLF.md},
\path{estimator/OBSERVER_ISS_THEORY.md}, \path{estimator/MODEL_AIDED_ESTIMATION.md},
\path{estimator/nrmmControllerErrorBounds.m} (the published target set),
\path{controller/predictiveSafetyGeometry.m} (\texttt{targetErrorTube}),
\path{scripts/collisionThreatContract.m} (the declared bounds of the
closed-loop target contract), \path{README.md}, \path{report/README.md},
\path{paper/README.md}, the local Zotero library (206 items) for the added
references, and \path{reference/vehicle_collision_avoidance_literature_review.md}.

\section{The derivation added (Section IV-J, Theorem 2)}
Let $\boldsymbol\theta=(\boldsymbol p_C,\gamma_C,V_C,A_C,\kappa_C)$ be the
obstacle's state and constants, $\varphi_\tau$ the flow of the forecast
(constant $A$ and $\kappa$, braking to rest, $A$ zero at rest), and
$\mathcal Q_k$ the set the estimator publishes at sample $k$. The flow is a
semigroup. The information set
$\mathcal S_{k+1}=\varphi_{T_s}(\mathcal S_k)\cap\mathcal Q_{k+1}$ contains the
truth, and the forecast families $\mathcal F_k(t)=\varphi_{t-t_k}(\mathcal S_k)$
are nested, $\mathcal F_{k+1}(t)\subseteq\mathcal F_k(t)$ (Lemma~6). This is
the unknown-but-bounded recursion of set-membership estimation applied to the
constants; the published sets alone are not nested. A computable nested outer
approximation intersects the published intervals of $A_C$ and $\kappa_C$ over
the history and takes, at each absolute time, the smallest of the enclosures
(eq.\ \texttt{target\_tube}) of the individual published sets evaluated with the
intersected intervals.

On the ego side the innovation
$\boldsymbol\delta_k=\hat{\boldsymbol x}_{k+1}-F(\hat{\boldsymbol x}_k,\boldsymbol u_k)
=-\boldsymbol\epsilon_{k+1}+F(\hat{\boldsymbol x}_k+\boldsymbol\epsilon_k,\boldsymbol u_k)
-F(\hat{\boldsymbol x}_k,\boldsymbol u_k)+\boldsymbol w_k$ lies in a fixed set
$\bar\Delta$ determined by the uniform bound of the ISS comparison solution
(Corollary~1 now states $\boldsymbol z(t)\le\sup_\tau\boldsymbol z(\tau)<\infty$),
the Jacobian bound of $F$ and the plant defect. A tube is a family of sets
$\mathcal D_s$ over the constrained instants $s\in\{\tfrac12,1,\tfrac32,\dots\}$
with (T1) $\mathcal D_{1/2},\mathcal D_1$ containing the true deviation of the
applied hold (the implemented per-hold propagation) and (T2)
$\mathcal D_{s+1}\supseteq\mathcal D_s\oplus\Phi_s\bar\Delta\oplus\mathcal R_s$
(Chisci-type growth, with $\Phi_s$ the sensitivity along the candidate's law,
open loop or in feedback form). The robust problem imposes the ego rows on
$\boldsymbol z_{s|k}\oplus\mathcal D_s$ without slack, and the obstacle rows
against every member of $\mathcal F_k$ with slacks; $V_k^\star$ is the minimal
slack sum.

Theorem~2: under the hypotheses of Theorem~1 along the loop, a plant equal to
$F$ up to a bounded defect, an obstacle on the constant-parameter model, a
tube with (T1)--(T2) and Assumption~3, feasibility of the robust problem is
kept at the next sample, $V^\star_{k+1}\le V^\star_k-\xi^\star_{0|k}$, the
applied hold separates the true vehicles by $d_{\rm m}$ less its slack, the
first-hold deficits are summable, and the zero set of $V^\star$ is invariant.
The proof is Huang's shifted-solution argument with the candidate's state at
instant $s$ contained in the old instant $s+1$'s tube (by (T2)) and the new
forecast family contained in the old one (by Lemma~6); the encounter end and
the parked check are monotone for the same reason. The theorem reduces to
Proposition~2 with exact information. A paragraph states what the
implementation realizes: (T1) only; forecasts from the centers of
$\mathcal Q_k$; no intersection. The two earlier measurements (open-loop
full-horizon ego tube infeasible; full-horizon target set 31 of 84 with one
collision from a shared slack) are read as the tube growing without feedback
and the width of the published family.

\section{Numbers}
No simulation or test was run. Every number of the results section is that of
the previous revision. One illustrative calculation was added: the position
radius of the forecast family under eq.\ \texttt{target\_tube} with the median
published radii of the estimated-state campaign (course 0.20~rad, speed
1.4~m/s, position 0.064~m) and the declared bounds of the closed-loop target
contract on the constants ($|A_C|\le1.1$~m/s$^2$, $|\kappa_C|\le\sin(0.055)/1.6
=0.034$~m$^{-1}$, from \path{scripts/collisionThreatContract.m} over the seven
scenarios widened by 10\%), for an obstacle at 8~m/s on a straight course:
4.7~m one second ahead, 12.7~m at two seconds, 23.9~m at three seconds
(\path{MANUSCRIPT_NARRATIVE_REWRITE_20261011/forecastFamilyWidth.py}, output in
\path{forecastFamilyWidth.txt}). The manuscript quotes ``about 5~m one second
ahead and more than 20~m three seconds ahead''.

\section{Manuscript changes}
\begin{itemize}
\item Abstract and introduction rewritten around the narrative; the
  contributions are the estimator, the PCBF controller, the transfer of the
  guarantee (Section IV-J) and the closed-loop evaluation.
\item Section II (Problem, Models and System): the task, the vehicle models,
  the architecture with the obstacle perception pipeline as a short part of
  it, and the safety argument whose link (v) now cites Theorem~2 as well as
  Proposition~2. The former Section III (Perception) is gone: its
  obstacle-perception subsection is condensed into Section II-C, and the
  road-curb extraction subsection, whose code was removed from the repository
  on October 6, 2026 and whose figure cannot be regenerated, is removed
  together with the figure.
\item Section III (Estimator): prose rewritten; the assumptions, Lemmas~1--4,
  Theorem~1 and the gain programs unchanged; Corollary~1 extended by the
  uniform bound; the published ego and target sets named ($\mathcal E_k$,
  $\mathcal Q_k$).
\item Section IV (Control): prose rewritten; Lemma~5, Assumption~3,
  Propositions~1--2 and every equation unchanged; the soundness of the
  tightening stated after eq.\ \texttt{collision\_tightening}; Section IV-I
  retitled ``Recursive feasibility with exact information''; new Section IV-J
  with Lemma~6 and Theorem~2; the scope subsection updated.
\item Section V (Results): numbers unchanged; two sentences connect the
  shifted-plan failures and the published target widths to Section IV-J.
\item Discussion and conclusion rewritten; the AI-assistance disclosure
  updated.
\item Seven references added from the Zotero library and the literature:
  Xu, Grizzle, Tabuada and Ames (2018); V\"or\"os et al.\ (2026); Alan et al.\
  (2023); Jeon, Zemouche and Rajamani (2019); Mayne, Rakovi\'c, Findeisen and
  Allg\"ower (2006); Schweppe (1968); Milanese and Vicino (1991).
\item Compiled with \texttt{latexmk}: 28 pages, no undefined reference or
  citation, no overfull line.
\end{itemize}

\section{Documentation changes}
\path{paper/README.md} rewritten for this revision; \path{report/README.md}
(this entry); the manuscript sentence of \path{CLAUDE.md} updated. No
executable statement and no controller or estimator document was changed.

\section{Open items for review}
\begin{itemize}
\item The proof of Theorem~2 and Lemma~6, and the semigroup argument for the
  flow with the rest rule, have not been independently reviewed. The theorem
  is conditional on the tube hypothesis (T2), which the implementation does
  not realize; the manuscript says so.
\item The abstract is 275 words; IEEE Transactions typically ask for at most
  250.
\item The bibliographic data of the three references that are not in the
  Zotero library (Mayne et al.\ 2006, Schweppe 1968, Milanese and Vicino
  1991) were entered from memory and should be checked against the
  publishers' records; the four taken from Zotero should be checked as well.
\item Whether the removal of the road-curb subsection is what the user wants;
  the architecture figure still shows the road-curb branch, labeled unused in
  the caption.
\item The estimator-alone results (Sections V-A and V-B) were not rerun.
\item The remarks on the implemented approximation (per-hold tightening,
  forecast centers) restate \path{controller/PCBF_CLF_ARCHITECTURE.md} and
  \path{controller/TERMINAL_SAFE_SET.md}, Section~9; the theory document
  \path{controller/OBSERVER_ROBUST_PCBF_THEORY.tex}, which predates the
  braking-continuation terminal set and claims no recursive feasibility, was
  not updated to the new result.
\end{itemize}
\end{document}
